\chapter{اللقاء الثاني والثلاثون: آيات الطلاق والعضل وبدايات أحكام الرضاع}

\section{مقدمة}
السلام عليكم ورحمة الله وبركاته. الحمد لله رب العالمين، وأشهد أن لا إله إلا الله وحده لا شريك له، وأشهد أن محمداً عبده ورسوله، صلى الله عليه وعلى آله وصحبه وسلم.

افتتح الشيخ هذا اللقاء بتلخيص جامع لما مضى في آيات الطلاق والعدة؛ لأن هذا الباب تتشعب مسائله، ويكثر فيه الاضطراب إن لم تستقر صورة المسألة في ذهن الطالب. ثم انتقل إلى بقيّة الآيات من قوله تعالى: \begin{ayah}وَلَا تَتَّخِذُوا آيَاتِ اللَّهِ هُزُوًا\end{ayah} إلى قوله تعالى: \begin{ayah}وَالْوَالِدَاتُ يُرْضِعْنَ أَوْلَادَهُنَّ\end{ayah}.

وكان من أجمل ما في هذا المجلس أن الشيخ لم يكتف بتقرير الأحكام، بل ربطها كلَّها بأصل كبير: أن الطلاق والرجعة والنكاح والرضاع أبواب تعبدية، لا مجالات هوى أو عبث أو مجرد أعراف اجتماعية.

\separator

\section{تلخيص الباب السابق: لماذا كان هذا التمهيد مهماً؟}
نبّه الشيخ في أول اللقاء إلى أن الكلام على قوله تعالى: \begin{ayah}وَالْمُطَلَّقَاتُ يَتَرَبَّصْنَ بِأَنْفُسِهِنَّ ثَلَاثَةَ قُرُوءٍ\end{ayah} لا يحسن فهمه إلا بعد تمييز أنواع المطلقات:
\begin{itemize}
    \item المطلقة قبل الدخول، وهذه لا عدة عليها.
    \item والآيسة من المحيض، والتي لم تحض، فعدتهما بالأشهر.
    \item والحامل، فعدتها وضع الحمل.
    \item والمطلقة ذات الأقراء، وهي التي يجري فيها هذا الباب على وجهه المفصل.
\end{itemize}

وكان مقصود الشيخ من هذا التلخيص أن يتدرب الطالب على جمع النصوص في الباب الواحد، وألا يحمل حكماً نزل في صورة مخصوصة على جميع الصور من غير تمييز.

\begin{fawaid}[حسن تصور الصور قبل الخوض في الخلاف أصل لا يستغنى عنه]
كثير من الاضطراب في أبواب الأحكام سببه الخلط بين الصور المختلفة. فإذا حُرِّرت الأنواع أولاً، سهل بعد ذلك فهم النصوص، وحسن تنزيل كل حكم على موضعه.
\end{fawaid}

\section{الطلاق في هدي الشرع لا في أهواء الناس}
أكد الشيخ أن الله لم يجعل الطلاق قراراً مطلقاً بيد الرجل يفعل به ما يشاء، بل جعله محكوماً بهدى الشرع: في وقته، وفي قصده، وفي عدده، وفي طريقة المراجعة بعده. فليس له أن يطلق كيف شاء، ولا أن يمسك بقصد الإضرار، ولا أن يكرر الطلاق والرجعة بلا نهاية.

ومن هذا الباب شدد على أن الطلاق لا ينفصل عن التقوى؛ فكما أن الرجل مأمور أن يتقي الله في إمساك زوجته، فهو مأمور أن يتقيه في فراقها، فيكون قصده الإصلاح، وفعله على وفق حدود الله، لا على وفق الغضب والهوى.

\separator

\section{تأويل (ولا تتخذوا آيات الله هزوا)}
\isnad{قال أبو جعفر:} نهى الله عباده أن يجعلوا أحكامه وآياته موضع لعب أو عبث. وذكر الشيخ هنا أن من صور هذا الباب ما كان يقع من بعض الناس في الطلاق والعتاق، فيوقع اللفظ ثم يزعم أنه لم يرد به شيئاً، أو أنه كان لاعباً.

والمعنى الذي قرره \rawi{الطبري} أن آيات الله وحدوده لا يصح أن تعامل معاملة المزاح الذي لا يترتب عليه أثر، لأن الشريعة أنزلت ليقام حكمها ويعظم أمرها، لا لتدخل في أبواب التلاعب والاحتيال.

\begin{fawaid}[أبواب الأسرة من أعظم ما يجب أن يعظم فيه حكم الله]
من أخطر صور الاستهانة أن يجري الطلاق أو الرجعة أو العتق أو سائر العقود الشرعية على ألسنة الناس على جهة العبث. فهذه الأبواب ليست هزلاً، لأنها تتعلق بالحقوق والأنساب والبيوت والذمم.
\end{fawaid}

\section{تأويل (واذكروا نعمة الله عليكم وما أنزل عليكم من الكتاب والحكمة)}
ربط \rawi{الطبري} هذا الأمر بتذكر نعمة الله بما ساقه لعباده من الهداية، والكتاب، والحكمة. فالنعمة هنا ليست مجرد الامتنان العام، بل منها أن الله لم يترك الناس سدى في هذه الأبواب، بل بيّن لهم حدودها وشرائعها.

ثم فسر \textit{الحكمة} هنا بما جرت به عادته في هذا الموضع: السنن التي بينها رسول الله صلى الله عليه وسلم للأمة، مما يوضح القرآن ويفصل مجمله.

وهذا من أجمل ما في الباب: أن الله لم ينزل الحكم مجرداً، بل أنزل معه الوعظ، وربطَه بتذكر النعمة، حتى لا يظن المكلف أن الحدود تضييق مجرد، بل هي رحمة وهداية وصيانة للنفوس والبيوت.

\begin{fawaid}[من النعمة أن الله بيّن الأحكام ولم يترك الناس لأهوائهم]
إذا استحضر العبد أن تفاصيل الطلاق والرجعة والنكاح والرضاع من نعمة الله عليه، سهل عليه التسليم لها، ورأى في الشرع رحمة وحفظاً، لا قيداً مجرداً.
\end{fawaid}

\section{تأويل (يعظكم به واتقوا الله واعلموا أن الله بكل شيء عليم)}
بيّن \rawi{الطبري} أن هذا وعظ من الله للمؤمنين: أن يتقوا ربهم في هذه الأحكام، وأن يعلموا أنه مطلع على ظاهر أعمالهم وخفي نياتهم. ولذلك أطال الشيخ في التنبيه على خطر أن يتظاهر الإنسان بالإصلاح وهو يريد الإفساد، أو أن يلبس ثوب الاحتساب وهو في الباطن طالب هوى أو دنيا.

فالآية تربي المكلف على أن الله لا ينظر إلى الصورة الظاهرة وحدها، بل يعلم ما في القلب، وسيجازي على ذلك. ومن هنا كانت التقوى في أبواب الأسرة آكد؛ لأن صور الصلاح قد تكون سهلة الادعاء، لكن الله يعلم المقاصد الباطنة.

\begin{fawaid}[العلم بإحاطة الله يطهر النية في أبواب المعاملات]
إذا أيقن العبد أن الله يعلم ما في نفسه، لم يجرؤ على أن يتخذ ألفاظ الشرع ستاراً لظلم الناس أو أكل الحقوق أو الإضرار بالأهل، لأن العلم الإلهي يفضح التزين الكاذب.
\end{fawaid}

\separator

\section{تأويل (وإذا طلقتم النساء فبلغن أجلهن فلا تعضلوهن أن ينكحن أزواجهن)}
ذكر \rawi{الطبري} أن هذه الآية نزلت في منع الأولياء من التضييق على المرأة إذا انقضت عدتها، وأرادت أن تعود إلى زوجها بنكاح جديد، وكان التراضي بينهما قائماً على المعروف.

ونقل الروايات في سبب النزول، ثم رجح قاعدة أوسع من خصوص القصة: أن الآية دالة على تحريم العضل على كل ولي يمنع من له عليها ولاية من نكاح من رضيت به، إذا كان النكاح مشروعاً والرضا معتبراً.

وهنا تظهر قاعدة نفيسة كان الشيخ يؤكدها: الفرق بين \textit{النزول} و\textit{التنزل}. فقد تنزل الآية أولاً في واقعة بعينها، لكن حكمها لا يختص بتلك الحادثة، بل يتناول كل من كان على مثل حالها.

\begin{fawaid}[العبرة بعموم الدلالة لا بخصوص الحادثة وحدها]
معرفة سبب النزول تعين على الفهم، لكن الآية لا تحبس على صاحب القصة، بل تتنزل على كل من شاركه في المعنى الذي نزل الحكم لأجله.
\end{fawaid}

\section{استنباط الطبري: دلالة الآية على اعتبار الولي في النكاح}
استنبط \rawi{الطبري} من هذه الآية دلالة قوية على أن للولي حقاً معتبراً في عقد النكاح؛ لأنه لو كانت المرأة تملك أن تزوج نفسها مطلقاً، أو أن تختار أي أحد ليكون وليها من غير اعتبار، لما كان لنهي الولي عن العضل هذا المعنى الواضح.

فكون الله نهى الولي عن المنع والتضييق يدل على أن له أثراً في باب الإنكاح، وأن استعماله هذا الأثر في غير موضعه ظلم وعدوان.

\begin{fawaid}[قد يستنبط الحكم من لوازم الخطاب لا من منطوقه المباشر فقط]
من فقه \rawi{الطبري} أنه لا يقف عند ظاهر الجملة وحده، بل ينظر إلى ما يلزم منها. فالنهي عن العضل هنا أفاد ضمناً أن للولي مدخلاً معتبراً في النكاح، وإلا لكان النهي خالياً من الفائدة.
\end{fawaid}

\section{تأويل (ذلك يوعظ به من كان منكم يؤمن بالله واليوم الآخر)}
قرر \rawi{الطبري} أن هذا الوعظ خاص الأثر بمن صدق بالله واليوم الآخر؛ لأنه هو الذي يتعظ بالخطاب، ويرى في هذه الأحكام ديناً يلقى الله به، لا مجرد تدابير اجتماعية قابلة للتلاعب.

ومن هنا أعاد الشيخ تقرير أن الإيمان بالله واليوم الآخر ليس معنى تجريدياً، بل يظهر أثره في مثل هذه المواطن الدقيقة: حين يملك الإنسان أن يمنع، أو يضار، أو يحتال، ثم يمنعه خوف الله من ذلك.

\separator

\section{بدايات باب الرضاع: (والوالدات يرضعن أولادهن حولين كاملين)}
دخل المجلس في آخره إلى أول آية الرضاع، فذكر \rawi{الطبري} أن الله جعل تمام الرضاعة حولين كاملين لمن أراد إتمامها. ثم تعرض الشيخ لما قد يشتبه من الجمع بين هذه الآية وبين قوله تعالى: \begin{ayah}وَحَمْلُهُ وَفِصَالُهُ ثَلَاثُونَ شَهْرًا\end{ayah}.

وبيّن أن آية الثلاثين شهراً ليست حداً تعبدياً لازماً في كل فرد من أفراد الناس، بل خبرٌ عن حال من الأحوال التي تقع في الخلق. أما آية الحولين فوردت في سياق التشريع والتحديد لمن أراد تمام الرضاعة، فهي من باب الحكم الذي يتعلق بقدرة المكلف وفعله.

ومن هنا قرر قاعدة منهجية مهمة: أن أمر الله ونهيه إنما يتعلقان بما يقدر العباد على فعله أو تركه. أما ما لا سبيل لهم إلى التحكم فيه، كمدة الحمل من حيث أصل وقوعها، فليس من جنس ما يتعلق به التكليف على ذلك الوجه.

\begin{fawaid}[لا يتعلق التكليف إلا بما يملكه العبد فعلاً أو تركاً]
من القواعد المحكمة أن الأمر والنهي إنما يردان على ما يطيقه العبد ويدخل تحت قدرته. أما ما لا سبيل له إلى إيجاده أو دفعه، فليس محلاً للتعبد بالأمر والنهي على هذا الوجه.
\end{fawaid}

\separator

ختم الشيخ اللقاء بالتأكيد على أن هذه الأبواب كلها لا يضبطها إلا أصلان: تعظيم آيات الله، ومعرفة حدود الله في كل باب قبل العمل فيه. فمن جمع بين هذين الأصلين سلم من كثير من الفوضى في النكاح والطلاق والرجعة والرضاع.
